\documentclass[10pt,a4paper]{amsart}
\usepackage[margin=19mm]{geometry}
\usepackage{amsmath,amssymb,mathtools}
\usepackage[scr=rsfs,cal=euler]{mathalfa}
\usepackage{xfrac}
\usepackage[unicode,hidelinks,bookmarks=false]{hyperref}
\newcommand{\ie}{i.e.\ }
\newcommand{\cf}{cf.\ }
\newcommand{\resp}{respectively\ }
\newcommand{\Subsection}{\subsection}
\providecommand{\nobreakdash}{\nobreak\hbox{-}}
\setlength{\emergencystretch}{2em}
\multlinegap0pt
\usepackage{mathtools}
\usepackage{bm}
\usepackage{dsfont}
\usepackage{tikz}
\newtheorem{assumption}{Assumption}
\newtheorem{lemma}{Lemma}
\usepackage{cleveref}
\crefname{assumption}{Assumption}{Assumptions}
\crefname{lemma}{Lemma}{Lemmas}
\newcommand{\Li}{\mathcal{L}}
\newcommand{\N}{\mathbb N}
\newcommand{\R}{\mathbb R}
\newcommand{\X}{{\mathscr X}}
\newcommand{\adState}{p}
\newcommand{\addSwitching}{\mathbb{S}}
\newcommand{\cOper}{\calC}
\newcommand{\calC}{\mathcal{C}}
\newcommand{\calT}{\mathcal{T}}
\newcommand{\coerAH}{\eta_H}
\newcommand{\coerAV}{\eta_V}
\newcommand{\coerM}{\eta_m}
\newcommand{\contA}{\gamma_a}
\newcommand{\contB}{\gamma_b}
\newcommand{\contC}{\gamma_c}
\newcommand{\contM}{\gamma_m}
\newcommand{\ddt}{\tfrac{\mathrm{d}}{\mathrm{d}t}}
\newcommand{\desired}[1]{#1_{\mathrm{d}}}
\newcommand{\dt}{\,\mathrm{d}t}
\newcommand{\inpVar}{u}
\newcommand{\inpVarDes}{\desired{u}}
\newcommand{\inpVarDim}{\rho}
\newcommand{\nrModes}{L}
\newcommand{\nrSwitches}{N}
\newcommand{\outVar}{y}
\newcommand{\outVarDes}{\desired{y}}
\newcommand{\outVarDim}{p}
\newcommand{\state}{\theta}
\newcommand{\switch}{\sigma}
\newcommand{\timeInt}{\mathbb{T}}
\title{Certified model predictive control for switched evolution equations using Model Order Reduction}
\author{Michael Kartmann${}^{\dagger}$ \and Mattia Manucci${}^\star$ \and Benjamin Unger${}^\ddagger$ \and Stefan Volkwein${}^\dagger$}
\date{Source: arXiv:2412.12930v2 (CC BY 4.0)}
\begin{document}
\maketitle

\noindent\emph{Source and licence.} Certified model predictive control for switched evolution equations using Model Order Reduction. Version arXiv:2412.12930v2 (CC BY 4.0), distributed by arXiv under the Creative Commons Attribution 4.0 International licence, http://creativecommons.org/licenses/by/4.0/, as recorded in the arXiv licence metadata for this exact version and shown on the abstract page https://arxiv.org/abs/2412.12930v2. Full licence notice: \texttt{licenses/arxiv-2412.12930-CC-BY-4.0.txt}.\par\smallskip

\noindent\textit{Editorial scope.} Counted unit: the article's lemma lemma:S-IBP (source0.tex lines 3286--3294). The article's complete original proof of it is delivered with it (source0.tex lines 3296--3305, one \texttt{proof} environment of 10 lines). No mathematics is written, repaired or re-derived: the statement and the proof are the article's own lines, in the article's own notation, and no retained line is substituted or re-flowed.  Retained uncounted as the source-local context the proof needs: lines 400--408; lines 465--507.  Nothing was omitted from any retained span, so the package carries no omission record.  No retained line contains a citation, so the wrapper carries no bibliography and no bibliography entry.  No retained line contains a figure environment, an image-inclusion command or drawing code, so no image is delivered and the package declares no asset dependency.  Editorial formatting only, in addition: an AMS \texttt{amsart} preamble that restates the article's own declarations and macros used by the retained lines, namely the environments \texttt{assumption}, \texttt{lemma} on separate counters, together with the standard packages those lines need.  The retained environments print in this document as Assumption 1, Lemma 1 in order of appearance, whereas the article prints the same items with the numbering of its own full document.  The complete native source of the article as distributed in the arXiv e-print for this exact version is bundled unchanged as \texttt{sources/170270/source0.tex}.
\begin{equation}
    \label{eqn:weakForm} %\tag{$SS(\theta_0, t_0,T)$}
    \left\{\quad
    \begin{aligned}
        m_{\switch}(\ddt \state,\varphi) + a_{\switch}(\state,\varphi) &= b_{\switch}(\inpVar,\varphi) && \text{for all $\varphi\in V$, a.e. (almost everywhere) in }\timeInt_n,\\
       \state(t_n) &= \state_{n} && \text{in } H,\\
       \outVar &= \cOper_{\switch}\state && \text{a.e. in }\timeInt_n, % \text{ f.a.a. } t>0,
    \end{aligned}\right.
\end{equation}

\begin{assumption}
	\label{ass:switchedSystem}
	Consider the switched system~\eqref{eqn:weakForm}.
	\begin{enumerate}
		\item\label{ass:switchedSystem:switching} The external switching signal $\switch$ belongs to the set
			\begin{align*}
				\addSwitching \vcentcolon= \left\{\switch\colon (0,\infty)\to \{1,\ldots,\nrModes\} \,\left|\,\begin{aligned}&\switch \text{ is piecewise constant with}\\
				 &\text{locally finite number of jumps}\end{aligned}\right.\right\}.
			\end{align*}
            For a given switching signal $\switch\in\addSwitching$, we define
   \begin{equation}\label{eqn:swi:tim}
				\calT_n \vcentcolon= \left\{t\in \timeInt_n \,\left|\,\switch \text{ is discontinuous in t} \right.\right\}.
    \end{equation}
    For simplicity, we assume that $\{t_n,T_n\}\notin \calT_n$ and that there are $|\calT_n| = \nrSwitches-1\in \N$ switches. Note that $N$ depends on $n$, but we suppress this subscript for readability. Finally, we denote the switching times by $t_{n,1}<\ldots<t_{n,i}<\ldots<t_{n,\nrSwitches-1}$ and set $t_{n,0}\coloneqq t_n$ and $t_{n,N}\coloneqq T_n$.
 		\item\label{ass:switchedSystem:A} The bilinear forms $a_i\colon V\times V$ are uniformly continuous and satisfy a uniform G\r{a}rding's inequality, i.e., there exist constants $\coerAV,\contA>0$ and $\coerAH\geq 0$ such that
        	\begin{equation}
        		\label{eqn:properties:a}
        		a_i(\phi,\phi)\geq \coerAV\|\phi\|_V^2 - \coerAH\|\phi\|_H^2,\qquad
        		a_i(\phi,\varphi) \leq \contA \|\phi\|_V\|\varphi\|_V,
       		 \end{equation}
       		 for all $\phi,\varphi\in V$ and all $i=1,\ldots,\nrModes$.
		\item\label{ass:switchedSystem:M} The bilinear forms $m_i\colon H\times H\to \R$ are symmetric and define inner products in the pivot space $H$. In particular, there exist constants $\coerM,\contM>0$ such that 
			\begin{equation}\nonumber
				%\label{eqn:properties:m}
				m_i(\varphi,\varphi)\geq \coerM\,{\|\varphi\|}_H^2,\qquad
				m_i(\phi,\varphi) \leq \contM\,{\|\phi\|}_H {\|\varphi\|}_H
			\end{equation}			        
			for all  $\phi,\varphi\in H$ and all $i=1,\ldots,\nrModes$.
		\item\label{ass:switchedSystem:B} The bilinear forms $b_i\colon \R^{\inpVarDim}\times V$ are uniformly continuous, i.e., there exists a constant $\contB>0$ such that
          	\begin{equation}\nonumber
          		%\label{eqn:properties:b}
          		b_i(\bar{\inpVar},\varphi) \leq \contB\, {\|\bar{\inpVar}\|}_{\R^{\inpVarDim}}\|\varphi\|_{V}
          	\end{equation}
          	for all $\bar{\inpVar}\in\R^{\inpVarDim}$, $\varphi\in V$ and all $i=1,\ldots,\nrModes$.
		\item\label{ass:switchedSystem:C}The operators $\cOper_i\in\Li(H,\R^{\outVarDim})$ are uniformly bounded, i.e., there exists a constant $\contC>0$ such that
			\begin{equation}\nonumber
          		%\label{eqn:properties:c}
          		\|\cOper_i\phi\|_{\R^{\outVarDim}} \leq \contC \|\phi\|_H
          	\end{equation}
          	for all $\phi\in H$ and all $i=1,\ldots,\nrModes$.
        \item The data satisfy $\inpVarDes\in L^\infty((0,\infty),\R^{\inpVarDim})$, $\outVarDes\in L^\infty((0,\infty),\R^{\outVarDim})$, $y_{T_n}\in\R^{\outVarDim}$ and $\theta_n\in H$.
 	\end{enumerate}
\end{assumption}

\begin{lemma}[Integration by parts]\label{lemma:S-IBP}
    Let $\switch\in\addSwitching$ and $\state,\adState\in \X_{n,\switch}$. Let $\calT_n$ be the set of switching times as in \Cref{ass:switchedSystem}\,\ref{ass:switchedSystem:switching} and assume that the bilinear forms $m_i$ satisfy \Cref{ass:switchedSystem}\,\ref{ass:switchedSystem:M}. Then
    \begin{align}    \label{eqn:integrationByPartsM}
    \begin{aligned}
        \int_{\timeInt_n} m_{\switch}(\dot{\state},\adState)\dt &=  m_{\switch(T_n)}(\state(T_{n}),\adState(T_{n}))-m_{\switch(t_n)}(\state(t_{n}),\adState(t_{n}))\\
        &+\sum_{t_i\in \calT} \left(m_{\switch(t^-_i)}(\state(t^-_{i}),\adState(t^-_{i}))-m_{\switch(t^+_i)}(\state(t_{i}^+),\adState(t^+_{i}))\right) -\int_{\timeInt_n} m_{\switch}(\dot{\adState},\state)\dt .
        \end{aligned}
    \end{align}
\end{lemma}

\begin{proof}
   First, observe that if $m_{\switch}$ and the adjoint state are continuous, the sum in \eqref{eqn:integrationByPartsM} would be zero, thus reducing the statement to standard integration by parts. Splitting $\timeInt_n$ into its sub-intervals with constant mode and applying integration by parts on each sub-interval, we can write
    \begin{align*}
        \begin{aligned}
            \int_{\timeInt_n} m_{\switch}(\dot{\state},\adState)\dt\;=&\;\sum_{i=0}^{N-1}\int_{t_i}^{t_{i+1}} m_{\switch}(\dot{\state},\adState)\dt\\
            =&\;\sum_{i=0}^{N-1}\big( m_{\switch(t^-_{i+1})}(\adState(t^-_{i+1}),\state(t_{i+1}^-)) - m_{\switch(t^+_i)}(\adState(t^+_{i}),\state(t_{i}^+))\big)-\int_{\timeInt_n} m_{\switch}(\dot{\adState},\state)\dt
        \end{aligned}
    \end{align*}
   The claim follows by rearranging the sum and from the fact that $t_n$ and $T_n$ have no switching points (see \Cref{ass:switchedSystem}\,\ref{ass:switchedSystem:switching}).
\end{proof}

\end{document}
